\setlength{\abovedisplayskip}{5pt}
\setlength{\belowdisplayskip}{5pt}
\setlength{\abovedisplayshortskip}{3pt}
\setlength{\belowdisplayshortskip}{3pt}

\noindent
\begin{minipage}[t]{0.48\textwidth}
\begin{enumerate}[label=\textbf{\arabic*.}, leftmargin=1.8em, itemsep=7pt, parsep=0pt, topsep=0pt, start=14]

\item \textbf{(a)} Các tinh thể natri clorat rất dễ nuôi cấy dưới dạng hình lập phương bằng cách để dung dịch nước và natri clorat bay hơi chậm. Nếu $V$ là thể tích của khối lập phương đó với độ dài cạnh là $x$, hãy tính $dV/dx$ khi $x = 3\text{ mm}$ và giải thích ý nghĩa của nó.
\par\vspace{3pt}
\textbf{(b)} Chứng minh rằng tốc độ thay đổi của thể tích khối lập phương theo độ dài cạnh bằng một nửa diện tích bề mặt của khối lập phương đó. Hãy giải thích theo quan điểm hình học tại sao kết quả này lại đúng bằng cách lập luận tương tự như Bài tập 13(b).

\item \textbf{(a)} Tìm tốc độ thay đổi trung bình của diện tích hình tròn theo bán kính $r$ khi $r$ thay đổi từ
\par\vspace{2pt}
\hspace*{0.5em}(i) $2$ đến $3$ \qquad (ii) $2$ đến $2{,}5$ \qquad (iii) $2$ đến $2{,}1$
\par\vspace{3pt}
\textbf{(b)} Tìm tốc độ thay đổi tức thời khi $r = 2$.
\par\vspace{3pt}
\textbf{(c)} Chứng minh rằng tốc độ thay đổi của diện tích hình tròn theo bán kính (tại mọi $r$) bằng chu vi của đường tròn đó. Hãy thử giải thích theo hình học tại sao điều này lại đúng bằng cách vẽ một hình tròn có bán kính tăng thêm một lượng $\Delta r$. Bạn có thể xấp xỉ độ thay đổi diện tích $\Delta A$ thu được như thế nào nếu $\Delta r$ là nhỏ?

\item Một hòn đá được thả xuống một cái hồ, tạo ra một gợn sóng tròn lan ra ngoài với tốc độ $60\text{ cm/s}$. Tìm tốc độ tăng của diện tích bên trong hình tròn sau \textbf{(a)} $1\text{ s}$, \textbf{(b)} $3\text{ s}$, và \textbf{(c)} $5\text{ s}$. Bạn có thể kết luận điều gì?

\item Một quả bóng bay hình cầu đang được bơm căng. Tìm tốc độ tăng của diện tích bề mặt ($S = 4\pi r^2$) theo bán kính $r$ khi $r$ là \textbf{(a)} $1\text{ ft}$, \textbf{(b)} $2\text{ ft}$, và \textbf{(c)} $3\text{ ft}$. Bạn có thể rút ra kết luận gì?

\item \textbf{(a)} Thể tích của một tế bào hình cầu đang lớn lên là $V = \frac{4}{3}\pi r^3$, trong đó bán kính $r$ được đo bằng micrômét ($1\text{ }\mu\text{m} = 10^{-6}\text{ m}$). Tìm tốc độ thay đổi trung bình của $V$ theo $r$ khi $r$ thay đổi từ
\par\vspace{2pt}
\hspace*{0.5em}(i) $5$ đến $8\text{ }\mu\text{m}$ \qquad (ii) $5$ đến $6\text{ }\mu\text{m}$ \qquad (iii) $5$ đến $5{,}1\text{ }\mu\text{m}$
\par\vspace{3pt}
\textbf{(b)} Tìm tốc độ thay đổi tức thời của $V$ theo $r$ khi $r = 5\text{ }\mu\text{m}$.
\par\vspace{3pt}
\textbf{(c)} Chứng minh rằng tốc độ thay đổi của thể tích hình cầu theo bán kính bằng diện tích bề mặt của nó. Giải thích theo hình học tại sao kết quả này đúng. Lập luận tương tự như Bài tập 15(c).

\item Khối lượng của một phần thanh kim loại nằm giữa đầu bên trái của nó và điểm cách đó $x$ mét về phía bên phải là $3x^2\text{ kg}$. Tìm mật độ dài (xem Ví dụ 2) khi $x$ là \textbf{(a)} $1\text{ m}$, \textbf{(b)} $2\text{ m}$, và \textbf{(c)} $3\text{ m}$. Mật độ cao nhất ở đâu? Thấp nhất ở đâu?

\item Nếu một bể nước hình trụ chứa $5000\text{ gallon}$, và nước thoát ra từ đáy bể trong $40\text{ phút}$, thì Định luật Torricelli cho biết thể tích $V$ của lượng nước còn lại trong bể sau $t$ phút là
\[
V = 5000\left(1 - \frac{1}{40}t\right)^2 \qquad 0 \leqslant t \leqslant 40
\]
Tìm tốc độ nước thoát ra khỏi bể sau \textbf{(a)} $5\text{ phút}$, \textbf{(b)} $10\text{ phút}$, \textbf{(c)} $20\text{ phút}$, và \textbf{(d)} $40\text{ phút}$. Vào thời điểm nào thì nước chảy ra nhanh nhất? Chậm nhất? Hãy tóm tắt các phát hiện của bạn.

\end{enumerate}
\end{minipage}\hfill
\begin{minipage}[t]{0.48\textwidth}
\begin{enumerate}[label=\textbf{\arabic*.}, leftmargin=1.8em, itemsep=7pt, parsep=0pt, topsep=0pt, start=21]

\item Lượng điện tích $Q$ tính bằng culông (C) đã đi qua một điểm trên dây dẫn cho đến thời điểm $t$ (tính bằng giây) được cho bởi công thức $Q(t) = t^3 - 2t^2 + 6t + 2$. Tìm cường độ dòng điện khi \textbf{(a)} $t = 0{,}5\text{ s}$ và \textbf{(b)} $t = 1\text{ s}$. (Xem Ví dụ 3. Đơn vị của dòng điện là ampe [$1\text{ A} = 1\text{ C/s}$].) Tại thời điểm nào thì cường độ dòng điện là nhỏ nhất?

\item Định luật vạn vật hấp dẫn của Newton phát biểu rằng độ lớn $F$ của lực do một vật thể có khối lượng $m$ tác dụng lên một vật thể có khối lượng $M$ là
\[
F = \frac{GmM}{r^2}
\]
trong đó $G$ là hằng số hấp dẫn và $r$ là khoảng cách giữa hai vật thể.
\par\vspace{3pt}
\textbf{(a)} Tìm $dF/dr$ và giải thích ý nghĩa của nó. Dấu trừ biểu thị điều gì?
\par\vspace{3pt}
\textbf{(b)} Giả sử biết rằng Trái Đất hút một vật thể với một lực giảm với tốc độ $2\text{ N/km}$ khi $r = 20\,000\text{ km}$. Lực này biến thiên nhanh như thế nào khi $r = 10\,000\text{ km}$?

\item Lực $F$ tác dụng lên một vật thể có khối lượng $m$ và vận tốc $v$ là tốc độ thay đổi của động lượng: $F = (d/dt)(mv)$. Nếu $m$ là hằng số, công thức này trở thành $F = ma$, trong đó $a = dv/dt$ là gia tốc. Nhưng trong thuyết tương đối, khối lượng của một hạt biến thiên theo $v$ như sau: $m = m_0/\sqrt{1 - v^2/c^2}$, trong đó $m_0$ là khối lượng của hạt khi nghỉ và $c$ là tốc độ ánh sáng. Hãy chứng minh rằng
\[
F = \frac{m_0 a}{(1 - v^2/c^2)^{3/2}}
\]

\item Một trong những nơi có thủy triều cao nhất trên thế giới xuất hiện tại Vịnh Fundy trên bờ biển Đại Tây Dương của Canada. Tại Mũi Hopewell, độ sâu của nước khi triều kiệt là khoảng $2{,}0\text{ m}$ và khi triều cường là khoảng $12{,}0\text{ m}$. Chu kỳ dao động tự nhiên lớn hơn 12 giờ một chút và vào một ngày trong tháng 6, triều cường đã xảy ra lúc 6:45 sáng. Điều này giúp giải thích mô hình sau đây cho độ sâu của nước $D$ (tính bằng mét) theo hàm của thời gian $t$ (tính bằng giờ sau nửa đêm) vào ngày đó:
\[
D(t) = 7 + 5\cos[0{,}503(t - 6{,}75)]
\]
Thủy triều dâng lên (hoặc hạ xuống) nhanh như thế nào vào các thời điểm sau?
\par\vspace{3pt}
\begin{tabular}{@{}p{0.45\linewidth}@{}p{0.45\linewidth}@{}}
\textbf{(a)} 3:00 sáng & \textbf{(b)} 6:00 sáng \\
\textbf{(c)} 9:00 sáng & \textbf{(d)} Buổi trưa
\end{tabular}

\item Định luật Boyle phát biểu rằng khi một khối khí bị nén ở nhiệt độ không đổi, tích của áp suất và thể tích là một hằng số: $PV = C$.
\par\vspace{3pt}
\textbf{(a)} Tìm tốc độ thay đổi của thể tích theo áp suất.
\par\vspace{3pt}
\textbf{(b)} Một khối khí được giữ trong một bình chứa ở áp suất thấp và được nén đều đặn ở nhiệt độ không đổi trong $10\text{ phút}$. Thể tích giảm nhanh hơn ở thời điểm bắt đầu hay ở thời điểm kết thúc $10\text{ phút}$ đó? Hãy giải thích.
\par\vspace{3pt}
\textbf{(c)} Chứng minh rằng độ nén đẳng nhiệt (xem Ví dụ 5) được cho bởi công thức $\beta = 1/P$.

\end{enumerate}
\end{minipage}
